\section{Calibration and complete-root comparisons}\label{sec:examples}
A detector is useful only after the possible differences and the allowed lower completions have been specified. The first calculation uses the supplied parity cochains to fix the four source classes of Section~\ref{sec:normalization}. The later bordism calculation distinguishes the remaining two product choices. Lower-dimensional examples test the same adopted equations and their complete gauge descent. None of these calculations uses a disputed real-space classification table.

\subsection{A single finite background fixes the source normalization}\label{subsec:finite-calibration}\label{sec:calibration04}
\subsubsection{The detecting background is a genuine integral lower solution}
Let
\begin{equation}\label{eq:detectG04}
G=C_2^{({\rho_1})}\times C_2^{(u)}\times C_2^{(v)},\qquad {s_1}=u+v,
\end{equation}
where $u,v,{\rho_1}$ are independent binary characters. The symbols $u,v$ are the two additional independent symmetry characters in this example, not gauge parameters. Define the signed integer CI two-cocycle
\begin{equation}\label{eq:detectk04}
{n_2}(g,h)=u(g)(-1)^{v(g)}v(h)\in Z^2(G,\Z_{s_1}),\qquad \bar n_2=uv.
\end{equation}
Take
\begin{equation}\label{eq:detecta04}
\boxed{\omega_2^{\mathbb R}=\tfrac12\bigl[({n_2})_{\rho_1+s_1}+({s_1}^2)_{\rho_1+s_1}\bigr],\quad
\check\omega_2^{\mathbb R}=\tfrac12({n_2})_{\rho_1+s_1},\quad B_3=-{\rho_1}_{\rho_1}\cup {n_2}_{s_1},\quad n_4=0.}
\end{equation}
The subscripts retain the coefficient systems. Direct substitution gives $\dd_{s_1}{n_2}=0$, $\dd_{\rho_1+s_1}\omega_2^{\mathbb R}\in\Z_{\rho_1+s_1}$, and
\begin{equation}
\mathcal O_5({n_2})=-B_3{n_2}-{\rho_1}{n_2}^2=0\quad\text{over the integers}.
\end{equation}
Thus this is a valid lower configuration, not an assignment of unrelated mod-two classes. On this background,
\begin{equation}
\bar B_3={\rho_1}{\bar n_2},\qquad \omega_2^{\SC}=2\omega_2^{\mathbb R}\bmod2=uv+(u+v)^2.
\end{equation}
The extension $\omega_2=\omega_2^{\mathbb R}\bmod1$ is half-valued. The restriction to the finite subgroup generated by the chosen symmetry lifts and fermion parity is therefore well-defined. Since $2\omega_2=0$ in $\R/\Z$ pointwise, it also defines the section $G\to G_b$ of the bosonic circle quotient used for the parity comparison. No uncontrolled ``large cyclic group'' approximation is involved.


\subsubsection{All charge-completion indeterminacies are included}
For this finite group,
\begin{equation}
H^*(BG;\F_2)=\F_2[{\rho_1},u,v].
\end{equation}
The mod-two reduction of the integral Bockstein for the sign character $\chi$ is
\begin{equation}\label{eq:twBock04}
D_\chi P=\Sq^1P+\chi P.
\end{equation}
The tensor product of the three cyclic periodic resolutions has differentials with entries $0,\pm2$. The integral K\"unneth calculation shows that its positive-degree cohomology with each of the sign systems used here has exponent two. Reduction of an integral class is therefore injective in the degrees in question. For the phase coefficient, the real-coefficient cohomology vanishes in positive degrees by finite-group averaging; hence the integral Bockstein detects phase classes. These statements are also checked directly in the finite resolution used below.

The possible parities of closed integer four-cochains of sign ${\rho_1}$ form $\ker D_{\rho_1}$ in degree four. A basis is
\begin{equation}\label{eq:closedCFbasis04}
\begin{gathered}
{\rho_1}v^3+v^4,\quad {\rho_1}uv^2+u^2v^2,\quad {\rho_1}u^2v+u^2v^2,\\
{\rho_1}u^3+u^4,\quad {\rho_1}^3v+{\rho_1}^2v^2,\quad {\rho_1}^3u+{\rho_1}^2u^2.
\end{gathered}
\end{equation}
A closed charge completion with parity $q_4$ changes the terminal phase class by
\begin{equation}\label{eq:CFimage04}
\tfrac12\bigl[\omega_2^{\SC}q_4+\Sq^2q_4\bigr].
\end{equation}
Its integral Bockstein reduction is $D_{s_1}(\omega_2^{\SC}q_4+\Sq^2q_4)$. Exact polynomial linear algebra gives rank four on the six-dimensional input in Eq.~\eqref{eq:closedCFbasis04}. The four $D_{s_1}(2\Theta_i)$ have rank four, and adjoining them raises the rank to eight. In particular none of the four deformations is hidden by an integer-charge shift.

For the independent parity computation we allow even more freedom: an arbitrary binary closed CF four-cochain, with no integer-lift requirement. The input dimension is fifteen, its phase-image rank is eight, and adjoining the four deformations gives rank twelve. The detecting cycles below annihilate this \emph{larger} image. This makes the comparison independent of which allowed parity CF completion the solver selects.


\subsubsection{Four exact detecting cycles}
Let $e_{a,b,c}$ denote the tensor product of the three standard cyclic periodic-resolution generators, in the order $({\rho_1},u,v)$. Its degree is $a+b+c$. With coefficient $\Z_{s_1}$, the three actions are $\chi=(+1,-1,-1)$, and
\begin{equation}\label{eq:periodicboundary04}
\partial_{s_1} e_{\boldsymbol m}
=\sum_{i:m_i>0}(-1)^{m_1+\cdots+m_{i-1}}
\begin{cases}
\chi_i-1,&m_i\text{ odd},\\
\chi_i+1,&m_i\text{ even}
\end{cases}
 e_{\boldsymbol m-\boldsymbol e_i}.
\end{equation}
Define degree-seven chains
\begin{align}\label{eq:detectorL04}
L_1&=e_{3,1,3}+e_{6,1,0},&
L_2&=e_{2,1,4},\nonumber\\
L_3&=e_{0,0,7}+e_{0,2,5},&
L_4&=e_{1,1,5}+e_{2,2,3}+e_{4,0,3}+e_{4,3,0},
\end{align}
and $Z_i=\partial_{s_1}L_i/2$. The division is exact, $\partial_{s_1}Z_i=0$, and $2Z_i$ is a boundary. Explicitly,
\begin{align}\label{eq:detectorZ04}
Z_1={}&e_{3,0,3}-e_{3,1,2}+e_{5,1,0}-e_{6,0,0},\nonumber\\
Z_2={}&e_{1,1,4}-e_{2,0,4},\nonumber\\
Z_3={}&-e_{0,0,6}-e_{0,2,4},\\
Z_4={}&e_{1,0,5}-e_{1,1,4}+e_{1,2,3}-e_{2,2,2}\nonumber\\
&+e_{3,0,3}-e_{4,0,2}+e_{3,3,0}-e_{4,2,0}.\nonumber
\end{align}
They obey
\begin{equation}\label{eq:detectmatrix04}
\boxed{\langle\Theta_j,Z_i\rangle=\tfrac12\delta_{ij}\pmod1,\qquad
\langle\text{every parity CF image},Z_i\rangle=0.}
\end{equation}
These are twisted ordinary group-homology cycles. We do not assert that each is the fundamental class of a manifold with a specified fermionic tangential structure. Their present role is to detect group-cochain obstruction classes, where they are sufficient. A future bordism comparison must additionally supply the structured cycles and the symmetry maps.

For reproducibility, map $e_{a,b,c}$ to the signed shuffle of $a$ copies of the ${\rho_1}$ generator, $b$ of the $u$ generator, and $c$ of the $v$ generator in the normalized bar resolution. The sign is the parity of inversions between distinct factors. In degree six this requires evaluating the $3^6=729$ ordered words once; the sums for all twenty-eight tensor cells are then available. It is an exact chain transfer, not random sampling.


\subsubsection{An independent target from the closed-Majorana parity theory}
The integer identity
\begin{equation}\label{eq:sheetremove04}
\dd_{s_1}(v)_{s_1}=-2{n_2}
\end{equation}
allows us to eliminate the parity sheet $n_2^{\SC}=2{n_2}$ by a full integer sheet gauge. Before that gauge, the lower map gives $\check n_3^{\SC}={\rho_1}{\bar n_2}$; the native Majorana field is $({\rho_1}+{s_1}){\bar n_2}$. The induced native Majorana field after Eq.~\eqref{eq:sheetremove04} is the explicit closed cochain
\begin{equation}\label{eq:gamma04}
\boxed{n_3^{\SC}=({\rho_1}+{s_1}){\bar n_2}+\omega_2^{\SC}\cup v+{s_1}\cup v^2\quad\text{over }\F_2.}
\end{equation}
Its cohomology class is $uv({\rho_1}+v)$, but the literal ordered expression in Eq.~\eqref{eq:gamma04}, rather than a silently reordered polynomial, is used in the evaluation.

Solve the remaining parity equation
\begin{equation}\label{eq:solveparity04}
\dd n_4^{\SC}=\Sq^2n_3^{\SC}+\omega_2^{\SC}n_3^{\SC}+{s_1}\Sq^1n_3^{\SC}.
\end{equation}
On normalized bar cochains of $C_2^3$, there are $7^4=2401$ CF unknowns and $7^5=16807$ equations. Deterministic exact elimination has rank $2092$ and a solution; we set the free variables to zero. Every equation is verified after solving. Changes of this solution lie in the larger parity CF image already annihilated in Eq.~\eqref{eq:detectmatrix04}.


Now evaluate the supplied closed-Majorana source~\cite[Sec.~V~G]{Project4}, without using a CI terminal formula to construct the target. In this example the names $n_3^{\SC}$ and $n_4^{\SC}$ refer to the fields \emph{after} the sheet-removing gauge. With the same lift and open-square rules as Section~\ref{sec:operations}, its entire phase is
\begin{align}\label{eq:closedparityeval04}
\Ob_6^{\SC}={}&\frac12\Big[
 \omega_2^{\SC}n_4^{\SC}+n_4^{\SC}\cup_2n_4^{\SC}
 +dn_4^{\SC}\cup_3n_4^{\SC}+dn_4^{\SC}\cup_4dn_4^{\SC}\Big]\nonumber\\
&+\frac12\Big[X_6(n_3^{\SC})+\zeta_{2,3}(\omega_2^{\SC},n_3^{\SC})\nonumber\\
&\qquad +(\Sq^2n_3^{\SC})\cup_4(\omega_2^{\SC}n_3^{\SC})
 +(\Sq^2n_3^{\SC})\cup_4(s_1\overline{\beta n_3^{\SC}})\nonumber\\
&\qquad +(\omega_2^{\SC}n_3^{\SC})\cup_4(s_1\overline{\beta n_3^{\SC}})
 +\zeta_{1,4}(s_1,\overline{\beta n_3^{\SC}})\nonumber\\
&\qquad +(\omega_2^{\SC}\cup_1s_1)\overline{\beta n_3^{\SC}}
 +s_1\Sq^2n_3^{\SC}\Big]\nonumber\\
&+\frac14\Big[\omega_2^{\SC}\beta n_3^{\SC}
 +(\beta n_3^{\SC})\cup_2(\beta n_3^{\SC})
 +s_1^2\big(\beta n_3^{\SC}+\overline{\beta n_3^{\SC}}\big)\Big].
\end{align}
The first two brackets are binary; the last retains signed integer arithmetic. The Cartan words are defined in Appendix~\ref{app:operations}, while the required Adem operation is
\begin{equation}\label{eq:X6-parity}
\begin{split}
X_6(x_3)={}&
 \big(\smile_{1213243142}+\smile_{1213431412}
      +\smile_{1232431421}+\smile_{1234314212}\big)(x_3,x_3,x_3,x_3)\\
 &+\overline{\beta x_3}\cup_2\overline{\beta x_3}.
\end{split}
\end{equation}
This is the fixed supplied representative, including its closed term. Twenty-four independent local degree-seven checks verify closure after the parity equation. The complete source is evaluated, not its individual nonclosed factors.

The sheet-removing gauge in Eq.~\eqref{eq:sheetremove04} changes a terminal source only by the differential of its full phase gauge. Thus its class may be computed after removing the sheet. Possible changes of the CF completion have already been removed from the detector. This is why the closed-Majorana sector is sufficient for this calibration; it is not a derivation of the entire $p+ip$ terminal phase anew.


\subsubsection{The full detecting matrix and a finite primitive}
The adopted formula \eqref{eq:4d-O6} and the independent target have the periods
\begin{equation}\label{eq:calperiod04}
\begin{array}{c|cccc}
&Z_1&Z_2&Z_3&Z_4\\\hline
\Ob_6&0&0&0&0\\
\Ob_6^{\SC}\text{ after sheet removal}&0&0&0&0\\
\Theta_1&1/2&0&0&0\\
\Theta_2&0&1/2&0&0\\
\Theta_3&0&0&1/2&0\\
\Theta_4&0&0&0&1/2
\end{array}
\end{equation}
Omitting the sheet-linear term $\rho_1\bar B_3\bar n_2/2$ from the adopted source would instead give $(0,0,0,1/2)$. Thus that term is necessary. Relative to the adopted formula, the coefficients of \emph{all four further} source deformations must be zero. The matrix detects operation classes modulo every allowed CF completion, rather than comparing four unrelated abstract group orders. Together with the fixed-model exhaustion theorem of Section~\ref{sec:normalization}, it fixes the finite-dyadic source normalization in that model.

There is a stronger finite check. For the deterministic parity CF completion, the transferred difference is exact without another CF shift. In the same small resolution its primitive is zero except for
\begin{equation}\label{eq:minprimitive04}
\begin{array}{c|rrrrrr}
(a,b,c)&(0,0,5)&(0,2,3)&(1,1,3)&(1,2,2)&(1,3,1)&(3,1,1)\\\hline
\widehat R_5^{\mathrm{min}}(e_{a,b,c})&1/4&5/48&1/8&5/8&1/8&3/8
\end{array}
\end{equation}
and all twenty-eight degree-six cells satisfy
\begin{equation}\label{eq:finite-primitive}
 d_{s_1}\widehat R_5^{\mathrm{min}}
 =\operatorname{sh}^*\Ob_6^{\SC}-\operatorname{sh}^*\Ob_6\pmod1.
\end{equation}
This finite primitive is not the missing universal circle-dependent map of \eqref{eq:remaining-R5}.

\subsection{A smaller mixed-background control in three dimensions}\label{subsec:mixed3}
Consider the finite quotient $G=C_2^{({\rho_1})}\times C_2^{({s_1})}$, with independent characters ${\rho_1},{s_1}$, CI root ${n_1}={s_1}$, and
\begin{equation}
\omega_2=\frac12{s_1}^2\in Z^2(G,\R/\Z_{{\rho_1}+{s_1}}).
\end{equation}
Here $\check\omega_2^{\mathbb R}=B_3=0$, so the integer source vanishes. The adopted finite-sum terminal source has zero image on the tensor minimal resolution of this quotient.

Now restrict the charge circle to $\Z_4^f$. Its bosonic quotient is another $C_2$, with generator $v$. The complete parity bosonic group is $C_2^3$, and
\begin{equation}
\check\omega_2^{\SC}=v^2+({\rho_1}+{s_1})v,\qquad \omega_2^{\SC}=\check\omega_2^{\SC}+{s_1}^2.
\end{equation}
The integer parity sheet $2{s_1}$ is exact. After it is removed with the corresponding parity gauge, the checked Majorana input is
\begin{equation}\label{eq:control-y}
\check n_2^{\SC}={\rho_1}{s_1}+\check\omega_2^{\SC}.
\end{equation}
This turns the comparison into an independent calculation with the project's \emph{closed-Majorana} source and phase. No CI terminal formula is used in that calculation.

The exact computation proceeds as follows. On the normalized bar complex of $C_2^3$, solve the binary CF equation for Eq.~\eqref{eq:control-y}. The linear system has rank $294$. Evaluate the fixed closed-Majorana phase from Ref.~\cite{Project3}, transfer it to the tensor minimal resolution by signed shuffles, and take the integral Bockstein. Since the relevant positive-degree integral cohomology with the sign local system has exponent two, its binary image detects the phase class. Repeat for the ten cubic binary CF cocycles, retaining their image in the terminal target.

The CF image has rank six. A permitted combination of the CF classes ${\rho_1}{s_1}v$ and ${s_1}^2v$ removes the terminal obstruction of this input. The raw phase and the selected CF completion are convention dependent; the vanishing of the remaining quotient class is the meaningful output. The minimal-resolution comparison itself was separately tested as a chain-map identity, for degrees one through five and all eight sign characters.


Suppose one changes a proposed $3+1$D CI phase by
\begin{equation}\label{eq:badclosed}
\Delta\Ob_5=\frac12{\rho_1}^3{s_1}\bar {n_1}.
\end{equation}
This is closed on the lower tower and is linear in the binary sheet. It therefore cannot be rejected merely by differentiating the source or its product equation. The elementary ${\rho_1}={s_1}$ root test alone is also not a sufficient detector of such universal changes.

For the mixed input above, Eq.~\eqref{eq:badclosed} is ${\rho_1}^3{s_1}^2/2$. Its integral Bockstein has the two nonzero minimal coordinates
\begin{equation}
({\rho_1},{s_1},v)\text{-degrees}\qquad (3,3,0),\quad(4,2,0).
\end{equation}
This vector is not in the CF image. It remains nonzero after every allowed CF completion change. Thus this particular alteration would disagree with the independently normalized parity calculation for the specified lower map.

This is a genuine normalization constraint, not a fit to a classification table. It excludes an actual closed change of the formula. It does \emph{not} enumerate or eliminate every possible closed natural operation, nor does a vanishing finite restriction by itself prove a continuous-group obstruction vanishes.


\subsection{Two-dimensional cyclic charge roots}\label{subsec:2d-examples}
For split unitary $\Uf\times C_m$, the completed roots are $x,y$ with
\begin{equation}
mx=-\binom m2y,\qquad my=0.
\end{equation}
The resulting relative group is $\Z_m^2$ for odd $m$ and $\Z_{m/2}\oplus\Z_{2m}$ for even $m$. The relevant self-twister can vanish in the chosen root representative even when the extension is nontrivial. The original phase completion and the integer gauge primitive that removes the accumulated charge must both be retained. The $3+1$D CI calculation below follows exactly this complete-state rule.

\label{sec:2d-bordism06}
For ${\rho_1}={s_1}=\omega_2=0$ and $G=C_2$, let $\alpha$ be its normalized character. The CF-only state
\begin{equation}
             x=(n_2,\nub_3)=(\alpha^2,\alpha^3/4)
\end{equation}
satisfies $\dd\nub_3=\alpha^4/2$. Its raw self-twister vanishes, but after stacking the integer charge $2\alpha^2$ is removed with $u_1=-\alpha$. The associated phase gauge is zero in this coordinate. Hence
\begin{equation}\label{eq:2d-root06}
                   2x\sim(0,\alpha^3/2),\qquad 4x\sim0.
\end{equation}
The bosonic sign is not removed by an incoming charge class, since $H^1(C_2;\Z)=0$. Thus the relative group computed here is $\Z_4$.

The independent comparison is
\begin{equation}
  \widetilde\Omega_3^{\operatorname{Spin}^c}(BC_2)
       \cong\Omega_2^{\operatorname{Pin}^c}\cong\Z_4,
\end{equation}
using the Smith equivalence of Ref.~\cite[Example~7.11]{SmithFiber} and the bordism value in Ref.~\cite[Theorem~9.87]{FreedHopkins}. This is a comparison of the \emph{relative} sector: uniform IQH and neutral $E_8$ response integers have been set to zero. It neither derives nor claims to normalize the entire nonzero-uniform-IQH product above. The script \texttt{low\_dimensional\_and\_incoming.py} independently checks Eq.~\eqref{eq:2d-root06} directly from the CF formulas, without calling a CI phase kernel.

\subsection{Three-dimensional AIII, AII, and AI}\label{subsec:3d-examples}
\label{sec:3d-bordism06}
These checks use the calibrated canonical convention, not a mixture of the short and canonical sources. Write $G=C_2$ with character $\alpha$. All displayed phases are cochain coordinates of completed states; only the phase of a state with all fermionic decorations removed is directly a pure bosonic response.

\paragraph{AIII: charge-conjugating antiunitarity.}
Set ${\rho_1}={s_1}=\alpha$ and $\omega_2=0$. The complete CI root and its reduced powers are
\begin{align}\label{eq:AIII-powers06}
 X&=(\alpha,0,-\alpha^4/16),\\
 2X&\sim(0,\alpha^3,\alpha^4/4),\\
 4X&\sim(0,0,\alpha^4/2),\qquad8X\sim0.
\end{align}
The first square uses the raw integer product $(2\alpha,\alpha^3)$ and $\Eb_4=\alpha^4/4$, followed by $u_0=1$ with $\widehat G_4=\alpha^4/8$. The following square includes the charge gauge. Thus $\Z_8$ is a complete-root relation, not the product of orders of three graded pieces.

The tangential structure is $\operatorname{Pin}^c$. The lift with antiunitary square one used here and the commonly used lift with square $(-1)^F$ differ by multiplying by an element of the charge circle; the circle is fixed by this antiunitarity ($(\rho_1+s_1)=0$). This change of section must not be confused with the distinct AI/AII extensions below. Independently,
\begin{equation}
           \Omega_4^{\operatorname{Pin}^c}\cong\Z_8\oplus\Z_2.
\end{equation}
The $\Z_8$ generator is detected by $\mathbb{RP}^4$ and the reduced eta invariant $\eta=\pm1/8$. Metlitski proves the character identity~\cite[Eq.~(74) and Sec.~VI]{Metlitski}
\begin{equation}\label{eq:AIII-eta06}
               e^{8\pi i\eta(M)}=(-1)^{\int_M w_1^4}.
\end{equation}
The relation $4X=b$ in Eq.~\eqref{eq:AIII-powers06} therefore matches a distinguished bordism character relation, not just an abstract group order. We have not evaluated the complete unsquared CI state on every $\operatorname{Pin}^c$ manifold. The independent neutral factor of the bordism group is not included in this CI-root calculation.

\paragraph{AII: charge-preserving time reversal with $T^2=(-1)^F$.}
Set ${\rho_1}=0$, ${s_1}=\alpha$, $\omega_2=\alpha^2/2$. Here $\check\omega_2^{\mathbb R}=B_3=0$, and the completed root
\begin{equation}
                         Y=(\alpha,0,0)
\end{equation}
satisfies $2Y\sim0$: its raw self-twister and the phase of the sheet-removing gauge $u_0=1$ are both zero. The bosonic sign $b=(0,0,\alpha^4/2)$ survives. The only relevant incoming integer CF class has
\begin{equation}
       \omega_2\alpha^2+\tfrac12\Sq^2\alpha^2
                      =\tfrac12\alpha^4+\tfrac12\alpha^4=0.
\end{equation}
The retained CI/charge/group-cohomological calculation is therefore $\Z_2\oplus\Z_2$. The full independent bordism group is
$\Omega_4^{\operatorname{Pin}^{\widetilde c+}}=(\Z_2)^3$~\cite{FreedHopkins}. The extra neutral response has not been constructed here; a mismatch of one factor is not evidence for changing the CI formula. Ref.~\cite[Proposition~6.27]{DebrayYeYu} gives independent generating manifolds $\mathbb{RP}^4,\mathbb{CP}^2,S^2\times S^2$ with their specified structures.

\paragraph{AI: charge-preserving time reversal with $T^2=1$.}
Set ${\rho_1}=0$, ${s_1}=\alpha$, $\omega_2=0$. The odd CI root has $Q_4(\alpha)=\alpha^4$ while an integer normalized degree-three cochain has differential $2n_3\alpha^4$. Thus the odd root is integrally obstructed. The charge layer has $H^3(C_2;\Z)=0$. The remaining group-cohomological bosonic sign is killed by the incoming degree-two CF class, whose phase image is $\alpha^4/2$.

Thus the computed CI/charge/group-cohomological sector is trivial. The full bordism group is $\Omega_4^{\operatorname{Pin}^{\widetilde c-}}=\Z_2$, generated by $\mathbb{CP}^2$ with its specified auxiliary bundle~\cite{FreedHopkins,DebrayYeYu}. This is again a neutral contribution outside the current sector, not a failure of the odd-CI obstruction. Keeping the integer obstruction and the incoming phase quotient is essential for this comparison.

\subsection{The decisive four-dimensional bordism test}\label{subsec:decisive-bordism}\label{sec:bordism06}
The deformation being tested is $\epsilon\Xi_5^{\mathrm{prod}}$ from Section~\ref{sec:normalization}. It does not alter the source, lower multiplication, or single-state gauge. We now compute a complete CI root whose square distinguishes the two values, and use an independent structured bordism calculation to select the adopted value $\epsilon=0$.
\subsubsection{The symmetry: unitary charge conjugation with square equal to parity}
Take $G=C_2=\langle c\rangle$, with normalized binary character $\alpha(c)=1$, and choose
\begin{equation}\label{eq:Hminus-background06}
 {\rho_1}=\alpha,\qquad {s_1}=0,\qquad
 \omega_2=\frac12\alpha^2,\qquad
 CQC^{-1}=-Q,\qquad C^2=(-1)^F.
\end{equation}
The notation $\alpha^j$ means the normalized inhomogeneous $j$-cochain whose only nonzero value is one on $(c,\ldots,c)$. Integer types are stated separately; a binary cochain is lifted before a fractional coefficient is applied. The fermionic group is
\begin{equation}
 G_f=\langle U(\theta),C\mid C U(\theta)C^{-1}=U(-\theta),\ C^2=U(\pi)\rangle
      \cong\operatorname{Pin}^{-}(2).
\end{equation}
This is \emph{unitary} charge conjugation, not ordinary antiunitary time reversal. Its bosonic quotient by fermion parity is $O(2)$. On oriented spacetime the tangential structure is
\begin{equation}\label{eq:Hminus06}
 H_-=(\operatorname{Spin}\times\operatorname{Pin}^{-}(2))/\langle(-1,-1)\rangle.
\end{equation}
We will derive the relevant bordism group rather than look up an ambiguously named ``Pin-c'' group.

For the explicit lift $\omega_2^{\mathbb R}=\alpha^2/2$,
\begin{equation}\label{eq:C2-lower06}
 \check\omega_2^{\mathbb R}=\frac12\alpha^2,\quad B_3=-\alpha^3\ \text{in }\Z_{\rho_1},
 \quad n_2=\alpha^2\ \text{in }\Z,
 \quad \mathcal O_5({n_2})=-B_3{n_2}-{\rho_1}{n_2}^2=0.
\end{equation}
Thus $n_4=0$ is a legitimate integer completion. The lower equation is checked over the integers, not just modulo two.


\subsubsection{A complete square, including its integer and phase gauges}
The exact evaluation of the calibrated six-cochain and five-cochain product gives
\begin{equation}\label{eq:C2-square-numbers06}
 \Ob_6(\alpha^2,0)(c^6)=\frac7{24},\qquad
 \Eb_5((\alpha^2,0),(\alpha^2,0))(c^5)=\frac12.
\end{equation}
Here $c^j$ abbreviates the tuple of $j$ copies of $c$. In the normalized $C_2$ complex,
\begin{equation}\label{eq:C2-differential06}
 \dd_\chi(\alpha^p)_\chi
     =\big[(-1)^{\chi(c)}+(-1)^{p+1}\big](\alpha^{p+1})_\chi.
\end{equation}
Consequently the complete root can be chosen as
\begin{equation}
 X=\left(\alpha^2,\ 0,\ \frac7{48}\alpha^5\right),
 \qquad \dd\nub_5=\Ob_6.
\end{equation}
The raw square has
\begin{equation}
      N_2=2\alpha^2,\qquad N_4=\alpha^4,
      \qquad\Vb_5=\frac{19}{24}\alpha^5
                         +\frac\epsilon2\alpha^5.
\end{equation}
The integer sheet gauge $u_1=-\alpha$, $v_3=0$ removes both integer outputs. Its \emph{calibrated} phase is
\begin{equation}
 \widehat G_5(2\alpha^2,\alpha^4;-\alpha,0)=\frac5{24}\alpha^5.
\end{equation}
Combining every term gives
\begin{equation}\label{eq:square-complete06}
 \boxed{\quad 2X\sim\left(0,0,\frac\epsilon2\alpha^5\right),\qquad
             4X\sim0.\quad}
\end{equation}
These rational numbers are values of the prescribed cochain coordinates. In particular $\nub_5(X)=7\alpha^5/48$ is \emph{not} by itself the physical partition function of $X$: the fermionic and CI factors must also be included. After the square and the complete gauge have removed those decorations, the remaining phase is a pure bosonic partition function and can be evaluated directly.

The calculation uses the full functions in \texttt{src/\_reference/checks\_v04/calibrated.py}; it is not a fit of an undetermined phase to the desired relation. The entire source equation is checked for the initial root, raw product, and gauged output. Because $C_2$ has one normalized bar generator in every degree, Eq.~\eqref{eq:C2-differential06} checks its whole normalized phase equation. The companion log records every intermediate value.

The bosonic class
\begin{equation}\label{eq:boson06}
              b=(0,0,\alpha^5/2)
\end{equation}
is not killed by the incoming integer CF class. That class has degree three, and its phase image is
\begin{equation}
  \omega_2\alpha^3+\frac12\Sq^2\alpha^3
       =\frac12\alpha^5+\frac12\binom32\alpha^5=0\pmod1.
\end{equation}
There is no incoming CI class in $H^1(C_2;\Z)$. More importantly, the manifold below proves $b$ remains nontrivial in the \emph{full} physical bordism theory, independently of the truncation of the decoration model.


\subsubsection{Smith reduction, with all tangential signs specified}
Let $E\to BO(2)$ be the tautological real plane and $L=\det E$, with ${\rho_1}=w_1(E)$. An $H_-$ structure can be described as a chosen spin structure on
\begin{equation}\label{eq:Hminus-spin06}
                         TM\oplus E\oplus L.
\end{equation}
Since $E\oplus L$ is oriented, this includes an orientation of $M$. Its degree-two condition is
\begin{equation}\label{eq:Hminus-w206}
                    w_2(TM)=w_2(E)+{\rho_1}^2.
\end{equation}
The chosen spin lift, not just this equality of characteristic classes, is part of the structure.

A transverse section of $L$ has zero locus $Y=\operatorname{PD}({\rho_1})$. Its normal line is $L|_Y$ and
\begin{equation}
                 TM|_Y=TY\oplus L|_Y.
\end{equation}
Thus the induced wall structure is a spin lift of $TY\oplus E\oplus2L$. The Whitney formula gives
\begin{equation}\label{eq:wall-classes06}
 w_1(TY)={\rho_1},\qquad w_2(TY)=w_2(E).
\end{equation}
This is the tangential $\operatorname{Pin}^{\widetilde c+}$ structure (class AII). One way to verify the sign without relying on a naming convention is that $E+2L$ and $3E$ have the same spin twist:
\begin{equation}
 w_1(E+2L)=w_1(3E)={\rho_1},\qquad
 w_2(E+2L)=w_2(3E)=w_2(E)+{\rho_1}^2.
\end{equation}
Their virtual difference is spin. Debray--Ye--Yu explicitly identify the $3E$ twist with $\operatorname{Pin}^{\widetilde c+}$ in Proposition~6.27 and its proof~\cite{DebrayYeYu}. Throughout this comparison the structures are on tangent bundles; changing to normal-bundle conventions interchanges the corresponding pin signs.

Apply the Smith fiber sequence~\cite[Theorem~5.1 and Corollary~5.8]{SmithFiber} to
\begin{equation}
 X=BO(2),\qquad V=E+L,\qquad W=L.
\end{equation}
The unit sphere bundle of $L$ is $BSO(2)\to BO(2)$; on it the structure in Eq.~\eqref{eq:Hminus-spin06} is $\operatorname{Spin}^c$. The resulting exact sequence contains
\begin{equation}\label{eq:smith-minus06}
 \Omega_5^{\operatorname{Spin}^c}\longrightarrow
 \Omega_5^{H_-}\xrightarrow{\mathrm{Sm}_L}
 \Omega_4^{\operatorname{Pin}^{\widetilde c+}}\longrightarrow
 \Omega_4^{\operatorname{Spin}^c}.
\end{equation}
The independent input is Freed--Hopkins, Theorem~9.87~\cite{FreedHopkins}:
\begin{equation}
 \Omega_5^{\operatorname{Spin}^c}=0,\quad
 \Omega_4^{\operatorname{Spin}^c}=\Z^2,\quad
 \Omega_4^{\operatorname{Pin}^{\widetilde c+}}=(\Z_2)^3.
\end{equation}
A homomorphism from a finite group to $\Z^2$ is zero. Exactness therefore proves
\begin{equation}\label{eq:bordism-minus06}
                   \boxed{\Omega_5^{H_-}\cong(\Z_2)^3.}
\end{equation}
Equation~\eqref{eq:bordism-minus06} is derived here from the cited theorem and the specified Smith twist; it is not represented as an entry copied from the literature.

There is no free five-dimensional response hidden in this comparison. The rational background ring is
\begin{equation*}
             H^*(BO(2);\Q)=\Q[p_1(E)].
\end{equation*}
The twist $E+L$ is oriented, and rational spin bordism is concentrated in degrees divisible by four. Hence $\Omega_6^{H_-}\otimes\Q=0$. The Anderson-dual exact sequence for invertible phases consequently gives
\begin{equation}
     \operatorname{TP}_5(H_-)
       \cong\operatorname{Hom}(\Omega_5^{H_-},\R/\Z)
       \cong(\Z_2)^3.
\end{equation}
In particular \emph{every} physical phase with this symmetry has square one on every closed structured five-manifold. No choice of basis for the three generators changes that statement.


\subsubsection{An actual physical detector: \texorpdfstring{$\mathbb{RP}^5$}{RP5}}
Take $M=\mathbb{RP}^5$ and let $L$ be its tautological real line. Set $E=\mathbf1\oplus L$. Then ${\rho_1}=w_1(L)$ and
\begin{equation}
 w(TM)=(1+{\rho_1})^6,\qquad w_1(TM)=0,\qquad w_2(TM)={\rho_1}^2,
 \qquad w_2(E)=0.
\end{equation}
Thus Eq.~\eqref{eq:Hminus-w206} holds, and a spin lift of $TM+E+L$ exists. This is a legitimate $H_-$ background, including the nontrivial spin--charge relation; it is not an ordinary spin manifold with an arbitrarily inserted $C_2$ gauge field. Its Smith wall is $\mathbb{RP}^4$, carrying the induced $\operatorname{Pin}^{\widetilde c+}$ structure in Eq.~\eqref{eq:wall-classes06}.

The pure bosonic state in Eq.~\eqref{eq:boson06} has
\begin{equation}\label{eq:RP5sign06}
                  Z_b(\mathbb{RP}^5,E)
                     =(-1)^{\int_{\mathbb{RP}^5}{\rho_1}^5}=-1.
\end{equation}
It therefore cannot be trivialized by adding missing neutral-sector identifications. On the wall this is the familiar sign $(-1)^{\int_{\mathbb{RP}^4}w_1^4}$. Its nontriviality as a $\operatorname{Pin}^{\widetilde c+}$ invariant is also explicit in the proof of Ref.~\cite[Proposition~6.27]{DebrayYeYu}.

\begin{theorem}[The remaining fixed-coordinate product coefficient]\label{thm:epsilon06}
With the adopted source, the retained integer product and joint gauges, and the standard normalization of a pure bosonic phase, physical compatibility with $H_-$ bordism requires
\begin{equation}\label{eq:epsilon-selected06}
                          \boxed{\epsilon=0.}
\end{equation}
\end{theorem}
\begin{proof}
By Eq.~\eqref{eq:square-complete06}, the two candidates give
$Z_{2X}(\mathbb{RP}^5,E)=(-1)^\epsilon$. Equation~\eqref{eq:bordism-minus06} requires this value to be $+1$ for every physical character. Equation~\eqref{eq:RP5sign06} rules out declaring the extra boson zero. Thus $\epsilon=1$ is incompatible with physical bordism, and $\epsilon=0$ is the only remaining candidate.
\end{proof}

This fixes a coefficient \emph{relative to the specified coordinates}. A paired change of source, product and phase can move closed terms between formulas. The theorem is not a claim that a certain monomial must be absent in every gauge convention. It also does not require knowing the partition function of the unsquared root $X$ on $\mathbb{RP}^5$; the exact square has removed all nonbosonic decorations.


\subsubsection{A second four-dimensional test: commuting \texorpdfstring{$C_2$}{C2}}
Set ${\rho_1}={s_1}=\omega_2=0$ on $G=C_2$. For the completed root
\begin{equation}
       X_0=(\alpha^2,0,\alpha^5/12)
\end{equation}
the full calibrated calculation gives
\begin{equation}\label{eq:commuting-powers06}
 \boxed{2X_0\sim(0,\alpha^4,\alpha^5/2),\qquad
        4X_0\sim(0,0,\alpha^5/2),\qquad8X_0\sim0.}
\end{equation}
For example $\Ob_6(X_0)=\alpha^6/6$, the first self-twister is zero, and the sheet-removing phase gauge is $\alpha^5/3$. These produce the first relation without guessing a group extension. The next square requires the integer CF gauge that removes $2\alpha^4$; it is included in the supplied log.

The independent Smith equivalence
\begin{equation}
 MT\operatorname{Spin}^c\wedge BC_2\simeq\Sigma MT\operatorname{Pin}^c
\end{equation}
is stated in Ref.~\cite[Example~7.11]{SmithFiber}. Therefore
\begin{equation}\label{eq:commuting-bordism06}
 \widetilde\Omega_5^{\operatorname{Spin}^c}(BC_2)
        \cong\Omega_4^{\operatorname{Pin}^c}
        \cong\Z_8\oplus\Z_2.
\end{equation}
The cochain calculation reproduces the eightfold CI direction. The independent neutral direction has not been constructed by the current CI formulas; it must not be erased to claim that our $\Z_8$ equals the full group. Likewise the untilded five-dimensional phase classification also has intrinsic $U(1)^f$ response terms; Eq.~\eqref{eq:commuting-bordism06} is the relative symmetry-protected comparison.


\subsubsection{A sign control: \texorpdfstring{$C^2=1$}{C squared equals one} is a different problem}
Keeping ${\rho_1}=\alpha,{s_1}=0$ but changing $\omega_2$ to zero gives an odd-sheet equation $-2n_4=-\alpha^5$. It has no integer solution. The incoming degree-three charge class now has phase image $\alpha^5/2$ and removes the bosonic sign in the CI/charge/group-cohomological sector.

The corresponding structure is
$H_+=(\operatorname{Spin}\times\operatorname{Pin}^+(2))/\Z_2$,
with $w_2(TM)=w_2(E)$. It can be represented by the oriented spin twist $E+3L$. Its Smith wall has
\begin{equation}
 w_1(TY)={\rho_1},\qquad w_2(TY)+{\rho_1}^2=w_2(E),
\end{equation}
namely $\operatorname{Pin}^{\widetilde c-}$, not $\operatorname{Pin}^{\widetilde c+}$. The same exact-sequence argument gives
\begin{equation}
                   \Omega_5^{H_+}\cong\Omega_4^{\operatorname{Pin}^{\widetilde c-}}
                        \cong\Z_2.
\end{equation}
This remaining bordism direction is outside the computed CI/charge/group-cohomological sector. In particular $\mathbb{RP}^5$ with $E=1+L$ is \emph{not} an $H_+$ detector: its $w_2(TM)={\rho_1}^2$ fails the $H_+$ condition. Interchanging these two symmetries would reverse the logic of the calibration.
